\chapter{État de l'art et étude de l'existant}
\section{Problématique et méthode de comparaison}
Un éditeur collaboratif doit synchroniser les textes sans ralentir la saisie. Le projet ajoute un crédit quotidien, un nombre exact de téléchargements et un PDF actuel derrière un lien stable, y compris lors d'un pic de fréquentation.

L'état de l'art décrit les mécanismes ; l'étude de l'existant observe leur emploi. Google Docs sert de référence pour l'édition généraliste, Overleaf pour la collaboration autour de sources et de PDF. Yjs est étudié comme option, sans lui attribuer un usage chez ces produits.

La comparaison distingue trois niveaux de preuve : \textbf{fonction documentée}, \textbf{mécanisme technique publié} et \textbf{déploiement documenté}. 

\section{Mécanismes de collaboration}
\subsection{Transformation opérationnelle : OT}
La transformation opérationnelle (OT) adapte une modification aux changements intervenus depuis l'état connu de son auteur. Elle ajuste, par exemple, les positions d'une édition lorsque l'autre rédacteur a inséré du texte avant la zone visée. Google en décrit le principe en 2010 \cite{google-ot}.

Un protocole OT gère aussi l'ordre des révisions, les opérations en attente et leurs acquittements. Le protocole Google publié permet l'affichage local optimiste, puis réconcilie les saisies avec celles reçues du serveur \cite{google-protocole}. Une frappe affichée n'est donc pas nécessairement acceptée durablement.

\subsection{Types de données répliquées sans conflit : CRDT}
Les CRDT organisent les données pour fusionner des modifications concurrentes. Yjs fournit des types partagés indépendants du transport et de l'éditeur ; ses mises à jour sont commutatives, associatives et idempotentes. Les répliques convergent après réception de toutes les mises à jour \cite{yjs-intro,yjs-updates}.

La convergence ne remplace ni les autorisations ni la persistance ou le contrôle du crédit. Le choix OT/CRDT doit aussi respecter les règles d'acceptation des modifications.

\subsection{Transport, présence et historique}
Le transport livre les messages, mais ne résout pas leurs conflits. Overleaf documente OT pour les réconcilier et WebSocket pour notifier les clients. Présences et positions complètent le contenu partagé \cite{overleaf-collaboration}.

\section{Google Docs : fonctionnement publié et limites de connaissance}
\subsection{Fonctions pertinentes}
Google documente les rôles de partage, l'historique et la restauration, ainsi que l'enregistrement et le téléchargement en PDF. Ces fonctions servent de repères pour le projet \cite{google-partage,google-historique,google-export}.

Une limite doit être prise en compte : l'aide Google annonce un maximum de 100 appareils ou onglets pouvant modifier simultanément un même fichier. 


\subsection{Architecture du protocole décrite en 2010}
Dans la description de 2010, chaque client conserve ses opérations en attente et sa dernière révision synchronisée. Le serveur traite et enregistre les opérations, acquitte leur auteur et transmet les changements aux autres clients, qui les réconcilient localement \cite{google-protocole,google-ot}.

\umlfigure{google-protocole.png}{1}{Séquence UML simplifiée du protocole Google publié en 2010 ; illustration documentaire, 
sans prétention sur la pile actuelle.}{fig:google}


\section{Overleaf : collaboration et chaîne de production PDF}
\subsection{Mécanisme de collaboration et outils identifiés}
Overleaf documente OT pour réconcilier les changements côté serveur et WebSocket pour les transmettre aux clients \cite{overleaf-collaboration}.

Le dépôt officiel déclare Node.js, TypeScript et des dépendances liées à React et CodeMirror. Ce manifeste décrit le code publié, sans établir l'usage de chaque dépendance ni l'architecture du service commercial. OT et les bibliothèques d'interface répondent à des rôles distincts \cite{overleaf-code}.

\subsection{Architecture sur site documentée}
La documentation sur site sépare les services \texttt{web}, \texttt{real-time}, \texttt{document-updater}, \texttt{docstore}, \texttt{filestore} et \texttt{clsi}. Elle associe Redis aux mises à jour actives, MongoDB aux documents et Docker à l'exécution de ces composants \cite{overleaf-microservices}.

Server Pro documente Docker et Nginx ; sa compilation s'exécute dans des conteneurs. Server CE utilise des sous-processus, selon la documentation consultée \cite{overleaf-infra,overleaf-microservices}.

\umlfigure{overleaf-composants.png}{1}{Diagramme UML de composants : vue simplifiée des éditions Overleaf sur site, d'après leur documentation. Les dépendances expriment les échanges principaux.}{fig:overleaf}



\subsection{Persistance et limites de transposition}
La documentation de sauvegarde indique que Redis conserve temporairement des mises à jour avant leur transfert vers MongoDB et recommande AOF. Un cache seul ne prouve donc pas la conservation des écritures acceptées \cite{overleaf-backups}.

Overleaf distingue l'édition de la production PDF. Le sujet porte sur du texte simple : il n'impose ni LaTeX ni un moteur TeX à l'application. LaTeX sert uniquement à produire ce cahier.

\section{Synthèse comparative et enseignements}
\begin{table}[H]\small
\begin{tabularx}{\linewidth}{@{}L{31mm}YYY@{}}\toprule
Critère & Google Docs & Overleaf & Projet étudié\\\midrule
Collaboration & OT dans les articles historiques de 2010 & OT et WebSocket documentés & Conservation et convergence obligatoires ; mécanisme à choisir\\
Historique & Consultation et restauration documentées & Historique présent dans le code et les services publiés ; modalités non évaluées ici & Consultation et restauration obligatoires\\
Production PDF & Téléchargement / impression PDF documentés & Service de compilation identifié sur site & PDF actuel, complet et cohérent\\
Infrastructure connue & Pile actuelle complète non établie & Docker, Nginx, Redis, MongoDB et services sur site documentés & AWS, une région et S3 imposés\\
Plafonds métier & Garanties spécifiques non établies par les sources étudiées & Garanties spécifiques non établies par les sources étudiées & Téléchargements exacts et crédit quotidien imposés\\
Validation & Aucune extrapolation de charge retenue & Aucune extrapolation de charge retenue & Tests et preuves propres au projet requis\\\bottomrule
\end{tabularx}
\caption{Comparaison limitée aux faits documentés et aux besoins du sujet.}
\end{table}

La conception doit traiter séparément la synchronisation, l'acceptation, la production d'un PDF cohérent et les quotas métier. La convergence ne garantit pas le nombre exact de téléchargements, pas plus que le stockage dans S3.

